\section{Complete assignments and their load--cost frontier}\label{sec:heterogeneous}
The relevant projection keeps operating cost as well as load. A whole
interval of early loads can be physical while the cheapest vehicle count
changes discontinuously across it. We characterize this joint frontier
first with fixed shared resources, then with a separate adequate-resource
specialization.

There are $n\ge1$ mandatory depot-to-depot services $A_i$ in $[0,1]$
and $n$ services $B_j$ in $[2,3]$. Batteries have capacity $B>0$ and
$0<a_i,b_j\le B$. A fixed declared bipartite graph $\Gamma$ specifies
which A-to-B shared-bus assignments are allowed. Every allowed transition
has zero movement time, energy and cost. Every service retains its standalone
bus option and its original depot charging access. At most $2n$ buses are used, each at fee $f>0$,
starting full and returning full by hour 4. One shared early connector
has power $K\ge0$ in $[1,2]$; one shared terminal connector has power
$H\ge0$ in $[3,4]$. Each bus accepts the corresponding connector power.
There is no other charging, reserve, loss, taper, auxiliary draw or switching
time. Powers $K,H$ are fixed resources, not multiples of $n$.
The graph is an operational compatibility restriction, not a selected bank
of duties used to approximate a larger declared opportunity set. Write
$\nu$ for its maximum matching cardinality, including $\nu=0$.

Define
\[
 S=\sum_i a_i+\sum_j b_j,\quad L=\max\{0,S-H\},\quad
 U=\min\{\sum_i a_i,K\},\quad r_{ij}=[a_i+b_j-B]_+,
\]
\[
 r_m=\min_{\substack{M\subseteq\Gamma\\|M|=m}}\sum_{(i,j)\in M}r_{ij},\qquad
 r_0=0,\qquad c_m=f(2n-m),
\]
where $M$ ranges over allowed partial bipartite matchings and
$m=0,\ldots,\nu$. Every such cardinality exists by deleting edges from
a maximum matching. A matching of size $m$ uses $2n-m$ buses.

\begin{theorem}[Fixed shared-resource load--cost projection]\label{thm:sharedprojection}
The model is feasible exactly when $L\le U$. In that case put
\[
 m_{\max}=\max\{0\le m\le\nu:r_m\le U\},\qquad
 m_0=\max\{0\le m\le\nu:r_m\le L\}.
\]
Its complete physical projection is
\begin{equation}\label{eq:sharedprojection}
 \bigcup_{m=0}^{m_{\max}}
 \{(E,S-E,c_m):\max\{L,r_m\}\le E\le U\}.
\end{equation}
The thresholds $r_m$ are nondecreasing and discretely convex. The
lowest operating-cost boundary $C(E)$ of the complete-plan hull starts
at $(L,c_{m_0})$, visits $(r_m,c_m)$ for $m_0<m\le m_{\max}$,
and is horizontal at $c_{m_{\max}}$ on
$[\max\{L,r_{m_{\max}}\},U]$.
\end{theorem}

For a fixed matching, early purchases satisfy $0\le x_i\le a_i$ and
$x_i\ge r_{ij}$ for a paired service. Their sum ranges over
$[R(M),\sum_i a_i]$. Serialize them at power $K$: the shared early
constraint is exactly $E\le K$. Every bus is present for terminal
charging, and individual refill demands sum to $S-E$. Serialization at
power $H$ is possible exactly when $S-E\le H$. Zero-power cases require
zero energy and no division. Appendix~\ref{app:heterogeneous} proves
both directions and the frontier geometry. Cardinalities above $\nu$ do
not exist; those with $r_m>U$ are resource-infeasible. Neither supplies
a fictitious upper-clipped endpoint.
The energy-only feasible set is simply $[L,U]$; it is the joint cost
frontier that matters for convexification.

\begin{corollary}[Support from the cost-labelled frontier]\label{cor:sharedsupport}
Suppose $L\le U$ and $G(E)=F(E,S-E)$ is finite, continuous and strictly
convex on $[L,U]$. Let $E^*$ uniquely minimize $C(E)+G(E)$. Then
\begin{equation}\label{eq:sharedsupport}
 D=\CH\quad\Longleftrightarrow\quad E^*\in
 \{L\}\cup\{r_m:m_0<m\le m_{\max}\}
 \cup[\max\{L,r_{m_{\max}}\},U].
\end{equation}
The gap is positive precisely in the open intervals between adjacent
retained thresholds. Admissible marginal-price support additionally
requires Proposition~\ref{prop:support}'s supply and price-attainment
qualification.
\end{corollary}

The clipped endpoint $L$ and intermediate physical thresholds on collinear
edges must be retained. If $L=U$ or $m_0=m_{\max}$ the gap is zero.
At $f=0$ the zero-matching branch attains every load at zero operating
cost, so the positive-fee interior criterion would fail.
The empty matching remains available even with isolated vertices or no
allowed edges. Consequently the feasibility test is unchanged; when
$\nu=0$, the entire projection has operating cost $2nf$ and its gap is zero.

\begin{corollary}[Classical ordered matching reduction for the complete graph]\label{cor:sortedmatching}
If $\Gamma$ contains every A-to-B edge, then for increasingly sorted
energies $a_{(i)},b_{(i)}$,
\[
 r_m=\sum_{i=1}^m[a_{(i)}+b_{(m+1-i)}-B]_+.
\]
\end{corollary}
Monotonicity permits selecting the $m$ smallest energies in each batch;
a convex positive-part exchange inequality then favors opposite pairing.
The appendix proves this classical assignment reduction. No matching
algorithm novelty is claimed.
The shortcut does not extend to arbitrary allowed graphs. With A energies
$(5,15)$, B energies $(14,15)$ and battery capacity 20, allowing only the
two corresponding-index edges gives $r_2=10$; the complete graph's crossed
assignment gives $r_2=9$. The general theorem uses the graph's actual
minimum-cost matching values and never invents a forbidden edge.

The fixed allowed-graph extension has independent source-only proof review;
its sparse-graph implementation qualification has not been performed.
The complete-crossmatching specialization has a separate, prospectively
frozen fourteen-case exact implementation qualification. Raw
SOC enumeration and serial charging witnesses matched the sorted-threshold
formulas in all fourteen cases; all twenty-four negative controls and the
smoke anchor passed. The cases include a physically attained intermediate
collinear threshold and a supported clipped lower boundary, as well as
zero-power, zero-fee, infeasible and continuous-response cases. This is an
internal implementation check, separate from the eight-case adequate-resource
qualification below. Fixed chargers also impose throughput limits as service
demand grows; this is no claim of indefinite scalability with $n$.
Nonzero transition consumption would make total energy matching-dependent;
transition time or staggered charging windows would invalidate the common-
window serial argument. Inserting an arbitrary directed timetable graph
without addressing those changes is therefore outside this theorem.

\subsection{A separate adequate-resource specialization}
We next study a different hardware policy in which the complete assignment set has an exact
energy projection. There are $n\ge1$ mandatory depot-to-depot services
$A_i$ occupying $[0,1]$ and $n$ mandatory services $B_j$ occupying $[2,3]$.
Their consumptions are $a_i$ and $b_j$, with
\begin{equation}\label{eq:batchassumptions}
 0<a_i\le20,\qquad0<b_j\le20,\qquad a_i+b_j\le30
 \quad\text{for every }i,j.
\end{equation}
Every A-to-B assignment is allowed, without movement time, energy or cost.
At most $2n$ identical buses may be used. Each has capacity 20 kWh,
starts full and returns full by hour 4. Its operating fee is strictly
positive, $f>0$.

Early charging in $[1,2]$ uses $n$ connectors of 10 kW each and an
aggregate grid limit of $10n$ kW. Terminal charging in $[3,4]$ uses $n$
connectors of 30 kW each and an aggregate grid limit of $30n$ kW. Per-bus
acceptance matches those connector powers. There is no other charging,
reserve, loss, taper or switching time. The grid limits are part of the
assumption: connector counts alone would not establish the simultaneous
allocation used below.

Since the services in each batch are simultaneous and have positive duration,
a bus can cover at most one A and one B. Thus a complete assignment is a
partial bipartite matching $M$. A matched pair uses one bus; unmatched
services each use a separate bus. All crossmatchings and all cardinalities
$m=0,\ldots,n$ are retained.

Define
\begin{align}
 S&=\sum_i a_i+\sum_j b_j,&u_i&=\min\{a_i,10\},&U&=\sum_i u_i,\\
 r_{ij}&=[a_i+b_j-20]_+,&
 r_m&=\min_{|M|=m}\sum_{(i,j)\in M}r_{ij},&c_m&=f(2n-m),
\end{align}
where $r_0=0$. The edge cost $r_{ij}$ is the minimum early purchase
needed by a bus covering $A_i$ then $B_j$. The thresholds $r_m$ are
minimum-cost cardinality-matching values, not samples from a plan bank.

\begin{theorem}[Exact heterogeneous fleet projection]\label{thm:batchprojection}
Under \eqref{eq:batchassumptions} and the stated resources, including $f>0$,
the complete physical projection is exactly
\begin{equation}\label{eq:batchprojection}
 \bigcup_{m=0}^n\{(E,S-E,c_m):r_m\le E\le U\}.
\end{equation}
The thresholds are nondecreasing and discretely convex:
\begin{equation}\label{eq:discreteconvex}
 0\le r_1-r_0\le r_2-r_1\le\cdots\le r_n-r_{n-1}.
\end{equation}
Let $m_0=\max\{m:r_m=0\}$. The lowest operating-cost boundary of the
complete-plan hull is the affine interpolant through
$(r_m,c_m)$ for $m=m_0,\ldots,n$, followed by the horizontal segment
$c_n$ on $[r_n,U]$.
\end{theorem}

The physical part of the theorem is the battery and connector argument.
A bus that has served $A_i$ can buy any $x_i\le u_i$ early; pairing it
with $B_j$ adds the lower bound $x_i\ge r_{ij}$. There are exactly $n$
A buses, so their individual charge intervals can be allocated simultaneously.
Terminal replenishment fits by grouping each matched bus alone and each
unmatched A/B pair together on one connector. The discrete-convexity step
uses standard bipartite min-cost-flow integrality and convex LP value
functions. It is not a new matching theorem. The complete proof is in
Appendix~\ref{app:heterogeneous}.

Use the same scaled supply family as the later replication calculation:
\begin{equation}\label{eq:batchsupply}
 F_n(E,L)=4E+\frac{E^2+L^2}{10n},\qquad
 G(E)=F_n(E,S-E),\qquad E_0=\frac{S-20n}{2}.
\end{equation}
Strict convexity makes the hull minimizing early load $E^*$ unique.

\begin{theorem}[A complete support criterion for the two-batch class]
\label{thm:batchsupport}
For strictly positive fee $f>0$ under Theorem~\ref{thm:batchprojection},
\begin{equation}\label{eq:batchzero}
 D=\CH\quad\Longleftrightarrow\quad
 E^*\in\{r_m:m=m_0,\ldots,n\}\cup[r_n,U].
\end{equation}
The gap is strictly positive exactly when $E^*$ lies strictly inside an
adjacent threshold interval $(r_{j-1},r_j)$, $j>m_0$.
If $d_j=r_j-r_{j-1}>0$, an intermediate threshold $r_m$ minimizes the
hull objective exactly when
\begin{equation}\label{eq:batchvertex}
 \frac{f}{d_{m+1}}\le G'(r_m)\le\frac{f}{d_m},
 \qquad m_0<m<n.
\end{equation}
\end{theorem}

The criterion includes weak support at ties. A threshold is physical, but
an interior point between consecutive thresholds interpolates operating
costs that no fleet at that early load can attain. Collinear adjacent
segments do not erase the intervening physical threshold. If $f=0$,
the conclusion about strict interiors would fail: the zero-matching
branch realizes every early load at zero operating cost, so the gap is
always zero. Strict positivity of the fee is therefore substantive.

The values themselves require only matching thresholds and scalar convex
minimizations. With $\clamp(v;\ell,u)=\min\{u,\max\{\ell,v\}\}$,
\begin{equation}\label{eq:batchD}
 D=\min_{0\le m\le n}\{c_m+G(\clamp(E_0;r_m,U))\}.
\end{equation}
For each sloping hull edge $j>m_0$, evaluate its objective at
\begin{equation}\label{eq:batchh}
 h_j=\clamp\left(E_0+\frac{5nf}{2d_j};r_{j-1},r_j\right).
\end{equation}
Compare these values with the minimum on the final horizontal segment.
Their minimum is $\CH$. If $m_0=n$, that horizontal segment alone
remains and $D=\CH$.

\subsection{A heterogeneous full-crossmatching obstruction}
Take $n=2$, $f=7$, A energies $(5,15)$ and B energies $(14,15)$.
Every crossmatching is admitted, with early-requirement matrix
\[
 (r_{ij})=\begin{pmatrix}0&0\\9&10\end{pmatrix}.
\]
The crossed full matching requires 9 early kWh; the other full matching
requires 10. Thus $r_0=r_1=0$, $r_2=9$, $S=49$, and $U=15$.
The full physical projection is
\begin{center}
\begin{tabular}{ccc}
\toprule
Matched services $m$ & Early energy interval & Operating cost\\
\midrule
0 & $[0,15]$ & 28\\
1 & $[0,15]$ & 21\\
2 & $[9,15]$ & 14\\
\bottomrule
\end{tabular}
\end{center}
Here $G(E)=2401/20-9E/10+E^2/10$. The physical optimum is the
crossed two-bus plan at $E=9$. The hull joins $(0,21)$ to $(9,14)$
and minimizes at $E^*=151/18$, strictly inside that edge. Exactly,
\begin{equation}\label{eq:heteroexample}
 D=\frac{2681}{20},\qquad\CH=\frac{434201}{3240},\qquad
 \Delta=\frac{121}{3240}>0.
\end{equation}
At the physical optimum's marginal price $(49/10,4)$, its bill plus
operating cost is $2181/10$; the one-pair zero-early complete fleet costs
217, so its regret is $11/10$.

The hull endpoints have full-inventory witnesses. For zero early charge,
pair the 5-energy A with the 14-energy B and serve the other services
separately. Terminal demands 19, 15 and 15 fit two connectors by serially
grouping the last two. For early charge 9, use the crossed full matching:
the 15-energy A bus buys 9 early, and both buses buy 20 terminally.
Every used battery starts and finishes at 20. Mixture weights $11/162$
and $151/162$ give the hull point. These numbers are analytically derived
and independently checked in a source-only proof review. A separately
frozen eight-case exact qualification of this adequate-resource class
matched raw SOC enumeration against the threshold formulas in all eight
cases, including ties and continuous response sets; all eighteen negative
controls and the smoke anchor passed. This qualifies those implementations,
not the fixed shared-resource extension above or operational prevalence.

The same example also has a fixed shared-resource realization: take
$B=20$, one early connector of power $K=15$ and one terminal connector
of power $H=49$. Then $L=0$, $U=15$ and the same thresholds give the
identical entire joint load--cost projection by
Theorem~\ref{thm:sharedprojection}. Hence the displayed $D,\CH$, gap
and regret also hold for this different hardware. Individual allocation
sets differ: the 15-energy A bus may now buy more than 10 early kWh.
The theorem supplies new serial realizations at powers 15 and 49;
the earlier parallel-connector witness timings are not reused as shared
connector schedules. This conclusion is analytical and is not covered by
the eight-case adequate-resource implementation qualification.

Both classes contain meaningful service-assignment choices, but retain
two simultaneous batches and zero-time, zero-energy allowed transitions. The
shared-resource theorem uses fixed connector powers; only the separate
adequate-resource specialization scales charging resources with $n$.
Neither covers the public varied-time case. The following two-service model makes the
same physical/hull distinction transparent and shows how connector and
equipment changes can preserve or remove it.
